% 04_implementacion.tex — Implementación: algoritmos y lenguajes
\section{Implementación}
\label{sec:implementacion}

\subsection{Arquitectura general}
\label{sec:impl:arquitectura}

Se implementaron \textbf{seis} solvers: algoritmos BRUTE y SMART en
Python~3.11, C++17 y Java~17.
Todos comparten el mismo formato de instancia y la misma
interfaz de l\'inea de comandos:
\begin{lstlisting}[language=bash,basicstyle=\ttfamily\footnotesize]
<programa> <instancia.3dm> [--algo brute|smart] \
  [--time-limit <s>] [--seed <int>] [--output <path>]
\end{lstlisting}
La salida estándar incluye el matching seguido de una línea de estadísticas:
\begin{lstlisting}
k
i1
...
ik
# stats: time_ms=<float> nodes=<int> algo=<brute|smart> n=<n> m=<m>
\end{lstlisting}

\subsection{Pre-procesamiento}
\label{sec:impl:preproceso}

Tras leer la instancia, cada solver construye índices invertidos:
$\texttt{inX}[x]$, $\texttt{inY}[y]$, $\texttt{inZ}[z]$: listas de índices
de tripletas que contienen el elemento correspondiente.
Este pre-procesamiento es $O(m)$ y se realiza una sola vez.

El estado de la búsqueda se almacena en tres \emph{bitmasks} de $n$ bits
(\texttt{usedX}, \texttt{usedY}, \texttt{usedZ}): el bit~$e$ vale~1 si el
elemento~$e$ ya está cubierto por el matching parcial.
Cuando $n\le 64$ (todos los experimentos) se usan enteros de 64~bits, lo que
permite test de conflicto y conteo de elementos libres (\texttt{popcount}) en
tiempo $O(1)$.

\subsection{Algoritmo BRUTE --- backtracking simple}
\label{sec:impl:brute}

\begin{algorithm}[htbp]
\DontPrintSemicolon
\caption{BRUTE: backtracking sobre las $m$ tripletas}
\label{alg:brute}
\KwIn{instancia $(n,T)$ con índices inversos}
\KwOut{mejor matching $M^*$ encontrado}
$chosen \leftarrow [\,]$;\quad $best \leftarrow [\,]$;\quad
$usedX = usedY = usedZ \leftarrow 0$\;
\SetKwFunction{recur}{recur}
\SetKwProg{Fn}{Función}{:}{}
\Fn{\recur{$i$}}{
  \If{$i = m$}{
    \If{$|chosen| > |best|$}{$best \leftarrow \text{copy}(chosen)$}
    \KwRet
  }
  $(x,y,z) \leftarrow T[i]$\;
  \If{bit $x$ libre en $usedX$ \textbf{and} bit $y$ libre en $usedY$
      \textbf{and} bit $z$ libre en $usedZ$}{
    Activar bits $x,y,z$;\quad $chosen.\text{append}(i)$\;
    \recur{$i+1$}\;
    Limpiar bits $x,y,z$;\quad $chosen.\text{pop}()$\;
  }
  \recur{$i+1$}\tcp*{rama ``saltar''}
}
\recur{$0$};\quad \Return{$best$}
\end{algorithm}

BRUTE explora el árbol binario de $2^m$ hojas potenciales, restringido a las
ramas factibles.  En el peor caso (instancias densas) el costo es $O(2^m)$.

\subsection{Algoritmo SMART --- MRV + bitmask + forward checking}
\label{sec:impl:smart}

SMART incorpora cuatro mejoras que reducen drásticamente el árbol de búsqueda:

\begin{enumerate}
  \item \textbf{Bitmask en $O(1)$.}  Conflictos y conteo de elementos libres
        se verifican con operaciones de bit.
  \item \textbf{Cota superior (upper bound).}  Al entrar en cada nodo se
        calcula
        $\texttt{upper} = |chosen| + \min(\texttt{freeX}, \texttt{freeY},
        \texttt{freeZ})$.
        Si $\texttt{upper}\le|best|$, el nodo se poda.
  \item \textbf{MRV (Minimum Remaining Values).}  Se elige el elemento libre
        $e$ con el menor número de tripletas \emph{disponibles} (no
        bloqueadas por \texttt{used*}).
        Si algún $e$ tiene 0 tripletas disponibles (elemento inalcanzable),
        se descarta en esta rama y se recursa sin él.
  \item \textbf{Forward checking implícito.}  Al fijar una tripleta
        $(x,y,z)$, los bits $x$, $y$, $z$ en \texttt{used*} se activan;
        cualquier tripleta que comparte alguno de estos bits queda
        automáticamente bloqueada para ramas descendientes.
\end{enumerate}

\begin{algorithm}[htbp]
\DontPrintSemicolon
\caption{SMART: backtracking con MRV + bitmask + poda}
\label{alg:smart}
\KwIn{instancia $(n,T)$ con índices inversos; bitmasks \texttt{usedX,Y,Z},
      \texttt{discardedX,Y,Z}}
\KwOut{mejor matching $M^*$ encontrado}
\SetKwFunction{recur}{recur}
\SetKwProg{Fn}{Función}{:}{}
\Fn{\recur{}}{
  $\textit{upper} \leftarrow |chosen| + \min(\texttt{freeX},\texttt{freeY},\texttt{freeZ})$\;
  \If{$\textit{upper} \le |best|$}{\KwRet\tcp*{poda por cota superior}}
  Elegir $(dim, e)$ libre y no descartado con mínimo \#tripletas disponibles\;
  \If{no hay tales $(dim,e)$}{
    \If{$|chosen|>|best|$}{$best\leftarrow\text{copy}(chosen)$}
    \KwRet
  }
  \If{$\#\text{disponibles}(e) = 0$}{
    Marcar $e$ como descartado;\quad \recur{};\quad Des-marcar\;
    \KwRet
  }
  \ForEach{tripleta $i$ disponible de $(dim,e)$}{
    $(x,y,z)\leftarrow T[i]$;\quad Activar bits;\quad $chosen.\text{append}(i)$\;
    \recur{}\;
    Limpiar bits;\quad $chosen.\text{pop}()$\;
  }
  Marcar $e$ como descartado;\quad \recur{};\quad Des-marcar\tcp*{opción ``no emparejar $e$''}
}
\end{algorithm}

\subsection{Invariante de determinismo}
\label{sec:impl:determinismo}

Los tres lenguajes procesan las tripletas y los índices en el mismo orden
(determinado por la semilla de entrada), lo que garantiza \textbf{igual número
de nodos explorados} entre C++, Java y Python para la misma instancia y semilla.
Esto fue verificado experimentalmente (Tabla~\ref{tab:correctness}).

\subsection{Bug en C++ con \texttt{--time-limit} bajo \texttt{-O2}}
\label{sec:impl:bug}

Se detectó que el flag \texttt{--time-limit} en la implementación C++ reporta
\texttt{timed\_out=1} de forma prematura cuando el binario se compila con
\texttt{-O2}.  La causa es una interacción entre la optimización del compilador
y el mecanismo de \texttt{std::condition\_variable} del hilo observador
(\textit{watcher thread}).  Con \texttt{-O0} el comportamiento es correcto.

\textbf{Workaround.}  El script de benchmarks usa \texttt{timeout TL\_s ./3dm}
a nivel shell para C++, en lugar del flag interno.  Esta decisión no afecta
la corrección de los resultados: el número de nodos explorados y el óptimo
reportado coinciden con los de Java y Python en todas las instancias
completadas.

\subsection{Complejidad y correctitud}
\label{sec:impl:correctitud}

\paragraph{BRUTE.}  Peor caso $O(2^m)$.  Por construcción, toda solución
devuelta es un 3-matching válido (las restricciones de \texttt{used*} se
verifican antes de cada inserción).

\paragraph{SMART.}  El árbol podado es exponencialmente menor en la
práctica, pero el peor caso sigue siendo exponencial (el problema es
NP-difícil).  La cota superior basada en $\min(\texttt{freeX},\ldots)$
es \emph{admisible}: nunca descarta la solución óptima.  La poda MRV
es un caso especial de la heurística de \emph{fail-first} estándar en CSP.

\paragraph{Correctitud cruzada.}  Los seis solvers fueron validados con
\texttt{tools/validator.py} y produjeron el mismo valor $\OPT(I)$ en todas
las instancias completadas (Sección~\ref{sec:experimentos}).
